
\subsection{弱无障碍与弗洛尔上同调的存在性}
\label{subsec:3-6-3}

由 \textbf{[Oh93]} 的附录和 \textbf{[Oh96II]} 可知：若
$L^{(0)}=L$、$L^{(1)}=\psi(L)$，其中 $\psi$ 是哈密顿微分同胚，则对子
$(L^{(0)},L^{(1)})$ 的弗洛尔上同调可以在一个关于 $L$ 的较弱条件下定义；
这个条件比一般对子 $(L^{(0)},L^{(1)})$ 所需的条件弱。这个事实对各种应用
都很重要。本小节讨论这一现象的代数对应物。另见下面的第 3.7.3 小节。

\begin{definition}{3.6.29}\label{def:3-6-29}
\begin{enumerate}[label=\textup{(\arabic*)},leftmargin=2.8em]
  \item 假设 $C$ 有幺元 $\mathbf e$。若元素
  $b\in(C[1])^0=C^1$ 满足
  $b\equiv0\pmod{\Lambda^+_{0,\mathrm{nov}}}$，并且对某个
  $c\in\Lambda_{0,\mathrm{nov}}^{+(0)}$ 有
  \begin{equation}
    \mathfrak m(e^b)=ce\mathbf e,
    \tag{3.6.30}\label{eq:3-6-30}
  \end{equation}
  就称 $b$ 是一个\emph{弱定界余链}\liftIndex{I111}，或称它是一个
  \emph{弱 Maurer--Cartan 解}。这里
  $\Lambda_{0,\mathrm{nov}}^{+(0)}$ 是
  $\Lambda^+_{0,\mathrm{nov}}$ 的零次部分。把所有弱定界余链构成的集合记作
  $\widehat{\mathcal M}_{\mathrm{weak}}(C)$。若
  $\widehat{\mathcal M}_{\mathrm{weak}}(C)\ne\varnothing$，就称 $C$
  \emph{弱无障碍}\liftIndex{I070}\liftIndex{I353}。

  \item 假设 $C$ 有同伦幺元。与定义 3.3.2 一样，我们有相伴的
  $A_\infty$ 代数
  \[
    (C^+,\mathfrak m^+)
    =\bigl(C\oplus\Lambda_{0,\mathrm{nov}}\mathbf e^+
      \oplus\Lambda_{0,\mathrm{nov}}\mathbf f,\mathfrak m^+\bigr),
  \]
  其幺元为 $\mathbf e^+$。若按上面 (1) 的意义有
  $\widehat{\mathcal M}_{\mathrm{weak}}(C^+)\ne\varnothing$，就称 $C$
  弱无障碍。
\end{enumerate}
这里提醒读者：符号 $e$ 是记录马斯洛夫指标的形式参数，其次数为 $2$。
\end{definition}

\begin{remark}{3.6.31}\label{rem:3-6-31}
当 $C$ 有严格幺元 $\mathbf e$ 时，也可以把它看作有同伦幺元。因此，
$\widehat{\mathcal M}_{\mathrm{weak}}(C^+)$ 与
$\widehat{\mathcal M}_{\mathrm{weak}}(C)$ 都可以考虑；二者由下式联系：
\[
  \widehat{\mathcal M}_{\mathrm{weak}}(C)\ni b
  \longmapsto b+ce\mathbf f
  \in\widehat{\mathcal M}_{\mathrm{weak}}(C^+).
\]
事实上，当 $C$ 有严格幺元时，可以选取 $(C^+,\mathfrak m^+)$，使定义
3.3.2 中的 $\mathfrak h_k=0$（$k\geq2$）。于是
\[
\begin{aligned}
  \mathfrak m^+(e^{b+ce\mathbf f})
  &=\mathfrak m(e^b)
    +\mathfrak m^+(e^b\,ce\mathbf f\,e^b)\\
  &=ce\mathbf e+ce\mathfrak m_1^+(\mathbf f)
   =ce\mathbf e+ce(\mathbf e^+-\mathbf e)
   =ce\mathbf e^+.
\end{aligned}
\]
反过来，设 $b^+\in C^+[1]^0$ 是
$\mathfrak m^+(e^{b^+})=ce\mathbf e^+$ 的一个解。由次数原因，可写成
\[
  b^+=b+c'e\mathbf f,
  \qquad c'\in\Lambda_{0,\mathrm{nov}}^{(0)}.
\]
（回忆 $\deg e^m=2m$ 且 $\deg\mathbf e=0$。）同样的计算给出
\[
  ce\mathbf e^+
  =\mathfrak m^+(e^{b+c'e\mathbf f})
  =\mathfrak m(e^b)+c'e\mathfrak m_1^+(\mathbf f)
  =\mathfrak m(e^b)+c'e(\mathbf e^+-\mathbf e).
\]
比较两边可得 $c=c'$ 且 $\mathfrak m(e^b)=ce\mathbf e$。

不过，当 $C$ 只有同伦幺元而没有严格幺元时，上述论证并不成立。（形式上，
或许可以在 $C$ 本身上考虑关于同伦幺元 $\mathbf e\in C$ 的方程
$\mathfrak m(e^b)=ce\mathbf e$ 的解集
$\widehat{\mathcal M}_{\mathrm{weak}}(C)$；但除非 $C$ 有严格幺元，
我们从不考虑 $C$ 本身上的这个方程。）当 $C$ 有同伦幺元时，所有关于
$C$ 的弱无障碍性的陈述都应理解为关于 $C^+$ 的陈述。由于 $C^+$ 与 $C$
同伦等价（见注 3.3.7 (3) 以及第 4.2 节的引理 4.2.55），我们在 $C$
只有同伦幺元时定义
\[
  \widehat{\mathcal M}_{\mathrm{weak}}(C)
  :=\widehat{\mathcal M}_{\mathrm{weak}}(C^+).
\]
需要提醒读者：若 $\mathbf e$ 不是严格幺元，下面的引理未必对 $C$ 本身上
形式方程 $\mathfrak m(e^b)=ce\mathbf e$ 的解成立。
\end{remark}

\begin{lemma}{3.6.32}\label{lem:3-6-32}
若 $b\in\widehat{\mathcal M}_{\mathrm{weak}}(C)$，则
\[
  \delta_{b,b}\circ\delta_{b,b}=0,
\]
其中 $\delta_{b,b}$ 如定义 3.6.6。
\end{lemma}

\begin{proof}
由 (3.2.19) 和~\eqref{eq:3-6-7}，
\[
\begin{aligned}
0
&=(\mathfrak m\circ\widehat d)(e^bxe^b)\\
&=\mathfrak m(e^b\delta_{b,b}(x)e^b)
  +\mathfrak m(\widehat d(e^b)xe^b)
  +(-1)^{\deg x+1}\mathfrak m(e^bx\widehat d(e^b))\\
&=(\delta_{b,b}\circ\delta_{b,b})(x)
  +ce\mathfrak m(e^b\mathbf e e^bxe^b)
  +(-1)^{\deg x+1}ce\mathfrak m(e^bxe^b\mathbf e e^b)\\
&=(\delta_{b,b}\circ\delta_{b,b})(x)
  +cex+(-1)^{\deg x+1}(-1)^{\deg x}cex\\
&=(\delta_{b,b}\circ\delta_{b,b})(x).
\end{aligned}
\]
\end{proof}

\begin{definition}{3.6.33}\label{def:3-6-33}
\begin{enumerate}[label=\textup{(\arabic*)},leftmargin=2.8em]
  \item 对 $b\in\widehat{\mathcal M}_{\mathrm{weak}}(C)$，定义
  \[
    HF(C,b)
    =\frac{\operatorname{Ker}(\delta_{b,b}\colon C^+\to C^+)}
    {\operatorname{Im}(\delta_{b,b}\colon C^+\to C^+)}
    \liftIndex{I106}.
  \]

  \item 由方程
  \[
    \mathfrak m(e^b)=\mathfrak{PO}(b)e\mathbf e
  \]
  定义函数
  \[
    \mathfrak{PO}\colon
    \widehat{\mathcal M}_{\mathrm{weak}}(C)
    \longrightarrow\Lambda_{0,\mathrm{nov}}^{+(0)}.
  \]
  称这个函数为\emph{势函数}\liftIndex{I079}。
\end{enumerate}
\end{definition}

\begin{remark}{3.6.34}\label{rem:3-6-34}
\begin{enumerate}[label=\textup{(\arabic*)},leftmargin=2.8em]
  \item 这里定义的势函数与物理学文献中引入的超势密切相关
  \textbf{[KKLM00, HoVa00]}。在某些情形下，它已经得到部分计算
  \textbf{[ChOh03]}。最简单的情形另见第 3.7.7 小节；第 3.7.3 小节将使用
  势函数。（另见 \textbf{[FOOO08I,08II]}。）

  \item 当 $C$ 有同伦幺元时，对
  \[
    b\in\widehat{\mathcal M}_{\mathrm{weak}}(C)
    =\widehat{\mathcal M}_{\mathrm{weak}}(C^+)
  \]
  定义了 $\delta_{b,b}\colon C^+\to C^+$；引理 3.6.32 在这种情形下也成立。
\end{enumerate}

若再假设
\begin{equation}
\left\{
\begin{aligned}
&\mathfrak m(x_0\otimes\mathbf f\otimes x_1\otimes\cdots
  \otimes x_{k-1}\otimes\mathbf f\otimes x_k)\in C
  &&(k\ne0,\ x_i\in BC[1]),\\
&\mathfrak m_1(\mathbf f)-\mathbf e^+\in C,
\end{aligned}
\right.
\tag{3.6.35}\label{eq:3-6-35}
\end{equation}
则 $\delta_{b,b}$ 把 $C$ 映到 $C$，因而其限制
$\delta_{b,b}\colon C\to C$ 有定义。此外，自然包含
\[
  (C,\delta_{b,b})\longrightarrow(C^+,\delta_{b,b})
\]
是上链同伦等价。
\end{remark}

\begin{lemma}{3.6.36}\label{lem:3-6-36}
设 $\mathfrak f\colon C_1\to C_2$ 是滤过 $A_\infty$ 同态。
\begin{enumerate}[label=\textup{(\arabic*)},leftmargin=2.8em]
  \item 对任意满足
  $b\in C_1[1]^0$、$b\equiv0\pmod{\Lambda^+_{0,\mathrm{nov}}}$ 的 $b$，定义
  \begin{equation}
    \mathfrak f_*(b):=\mathfrak f(e^b)
    =\mathfrak f_0(1)+\mathfrak f_1(b)+\mathfrak f_2(b,b)+\cdots.
    \tag{3.6.37}\label{eq:3-6-37}
  \end{equation}
  这里
  \[
    \mathfrak f(e^b)=\left(\sum_{k=0}^{\infty}\mathfrak f_k\right)(e^b).
  \]
  它诱导映射 $\mathfrak f_*\colon C_1[1]^0\to C_2[1]^0$，并且
  $\mathfrak f_*(b)\equiv0\pmod{\Lambda^+_{0,\mathrm{nov}}}$。

  \item 假设 $C_i$ 有幺元，并令
  $\mathfrak f\colon C_1\to C_2$ 是含幺滤过 $A_\infty$ 同态（见定义
  3.3.11）。则对任意
  $b\in\widehat{\mathcal M}_{\mathrm{weak}}(C_1)$，都有
  \[
    \mathfrak f_*(b)=\mathfrak f(e^b)
    \in\widehat{\mathcal M}_{\mathrm{weak}}(C_2).
  \]
  因此，$\mathfrak f$ 诱导映射
  \[
    \mathfrak f_*\colon
    \widehat{\mathcal M}_{\mathrm{weak}}(C_1)
    \longrightarrow\widehat{\mathcal M}_{\mathrm{weak}}(C_2).
  \]
  当 $C_i$ 无障碍时，在下面的证明中取 $c=0$，还得到映射
  $\mathfrak f_*\colon\widehat{\mathcal M}(C_1)\longrightarrow
  \widehat{\mathcal M}(C_2)$。

  \item 进一步，在 (2) 的情形下，令
  \[
    \mathfrak{PO}_i\colon
    \widehat{\mathcal M}_{\mathrm{weak}}(C_i)
    \longrightarrow\Lambda_{0,\mathrm{nov}}^{(0)}
  \]
  为势函数。则
  \[
    \mathfrak{PO}_2\circ\mathfrak f_*(b)=\mathfrak{PO}_1(b).
  \]
\end{enumerate}
\end{lemma}

\begin{proof}
(1) 因为 $\mathfrak f$ 保持移位次数，显然
$\mathfrak f_*(b)\in C_2[1]^0$。又因为
$\overline{\mathfrak f}_0=0$，$\mathfrak f$ 保持能量滤过（见
(3.2.27.1)），并且
$b\equiv0\pmod{\Lambda^+_{0,\mathrm{nov}}}$，所以
$\mathfrak f_*(b)\equiv0\pmod{\Lambda^+_{0,\mathrm{nov}}}$。

(2) 令 $\widehat{\mathfrak f}$ 是由 $\mathfrak f$ 诱导的余代数同态，见
(3.2.28)。应用引理 3.6.2，可以把
$\widehat{\mathfrak f}(e^b)$ 写成指数形式。事实上，可以验证
\begin{equation}
  \widehat{\mathfrak f}(e^b)=e^{\mathfrak f(e^b)}.
  \tag{3.6.38}\label{eq:3-6-38}
\end{equation}
直接计算得到
\[
\begin{aligned}
\widehat d\bigl(e^{\mathfrak f(e^b)}\bigr)
&=\widehat d\bigl(\widehat{\mathfrak f}(e^b)\bigr)
 =\widehat{\mathfrak f}\bigl(\widehat d(e^b)\bigr)\\
&=ce\widehat{\mathfrak f}(e^b\mathbf e e^b)
 =ce\,e^{\mathfrak f(e^b)}\mathbf e e^{\mathfrak f(e^b)}.
\end{aligned}
\]
另一方面，
\[
  \widehat d\bigl(e^{\mathfrak f(e^b)}\bigr)
  =e^{\mathfrak f(e^b)}
   \mathfrak m\bigl(e^{\mathfrak f(e^b)}\bigr)
   e^{\mathfrak f(e^b)}.
\]
比较两式便得到所需的
$\mathfrak m(e^{\mathfrak f(e^b)})=ce\mathbf e$。(3) 也由上述计算得出。
\end{proof}

当 $C$ 是同伦含幺滤过 $A_\infty$ 代数时，也可以用同样的方式在
\[
  \widehat{\mathcal M}_{\mathrm{weak}}(C)
  =\widehat{\mathcal M}_{\mathrm{weak}}(C^+)
\]
上定义势函数。

下面研究弱无障碍情形下定理 3.6.18 的对应版本。我们采用如下假设。

\begin{statement}{假设}{3.6.39}\label{assump:3-6-39}
还假设当 $k<0$ 时
$H^k(\overline C,\overline{\mathfrak m}_1)=0$；并且
$H^0(\overline C,\overline{\mathfrak m}_1)$ 由幺元的等价类
$[\mathbf e]$ 生成。
\end{statement}

令 $G\subset\mathbb R_{\geq0}\times2\mathbb Z$ 是满足条件 3.1.6 的
子幺半群，并假设 $C$ 是 $G$-有隙的。与前面一样，把它写成
\[
  G=\{(\lambda_i,\mu_i)\mid i=0,1,2,\ldots\},
  \qquad i<j\Longrightarrow\lambda_i\leq\lambda_j.
\]

\begin{definition}{3.6.40}\label{def:3-6-40}
若 $G$-有隙 $A_\infty$ 代数 $C$ 存在一个弱定界余链 $b$，可以写成
\[
  b=\sum_{(\lambda_i,\mu_i)\in G}
    T^{\lambda_i}e^{\mu_i/2}b_i,
  \qquad
  b_i\in\overline C[1]^{-\mu_i}=\overline C^{1-\mu_i},
\]
就称 $C$ \emph{在 $G$ 上弱无障碍}\liftIndex{I070}\liftIndex{I112}。
\end{definition}

现在置
\begin{equation}
  G_{\mu\leq0}=\{(\lambda,\mu)\in G\mid\mu\leq0\}.
  \tag{3.6.41}\label{eq:3-6-41}
\end{equation}
从给定的 $G$ 枚举中删去所有满足 $\mu_i>0$ 的 $i$，便得到一个非递减函数
$j\colon\mathbb Z_{\geq0}\to\mathbb Z_{\geq0}$，满足 $j(0)=0$，并且
\begin{equation}
  G_{\mu\leq0}
  =\{(\lambda_{j(i)},\mu_{j(i)})\mid i=0,1,2,\ldots\}.
  \tag{3.6.42}\label{eq:3-6-42}
\end{equation}

\begin{theorem}{3.6.43}\label{thm:3-6-43}
令 $(C,\mathfrak m)$ 是 $G$-有隙滤过 $A_\infty$ 代数。在假设 3.6.39
下，存在 $o_k(\mathrm{weak})$ 和 $b_i(\mathrm{weak})$，具有下列性质：

\smallskip
\begin{equation}
  [o_k(\mathrm{weak})]
  \in H^{2-\mu_k}(\overline C,\overline{\mathfrak m}_1),
  \qquad
  b_i(\mathrm{weak})\in\overline C^{1-\mu_i}.
  \tag{3.6.44.1}\label{eq:3-6-44-1}
\end{equation}

\noindent\textup{(3.6.44.2)}
若 $i<k$ 时的 $[o_i(\mathrm{weak})]$ 与 $b_i(\mathrm{weak})$ 都已定义，
并且 $\mu_k\leq0$，则闭上链 $o_k(\mathrm{weak})$ 有定义；
$[o_k(\mathrm{weak})]$ 依赖于 $i<k$ 时的
$[o_i(\mathrm{weak})]$ 与 $b_i(\mathrm{weak})$。

\noindent\textup{(3.6.44.3)}
当 $\mu_k\leq0$ 时，若 $[o_k(\mathrm{weak})]$ 已定义且为零，
则 $b_k(\mathrm{weak})$ 有定义。当 $\mu_k>0$ 时，若所有 $i<k$ 的
$b_i(\mathrm{weak})$ 都存在，则 $b_k(\mathrm{weak})$ 存在。

\noindent\textup{(3.6.44.4)}
$(C,\mathfrak m)$ 在 $G$ 上弱无障碍，当且仅当可以归纳地选取
$b_i(\mathrm{weak})$\translatornote{原书此处写未定义的 $b'_i$；本定理开头
只引入 $b_i(\mathrm{weak})$，而 (3.6.44.2)、(3.6.44.3) 也始终用这一
序列，故据此统一。}，使得 $\mu_k\leq0$ 时的所有
$[o_k(\mathrm{weak})]$ 都为零。
\end{theorem}

换言之，只要假设所有满足 $\mu_i\leq0$、$i<k$ 的闭上链
$o_i(\mathrm{weak})$ 都已定义，并且
$[o_i(\mathrm{weak})]=0$ 于
$H^*(\overline C,\overline{\mathfrak m}_1)$ 中，就可以定义障碍类
$[o_k(\mathrm{weak})]$\liftIndex{I108}\liftIndex{I349}。

$o_k(\mathrm{weak})$ 由与 $o_k$ 相同的公式（定义 3.6.23）给出，而
$b_k(\mathrm{weak})$ 由下式定义：
\begin{equation}
  -\mathfrak m_{1,0}\bigl(b_k(\mathrm{weak})\bigr)
  =
  \begin{cases}
    o_k(\mathrm{weak}),&\mu_k\leq0,\\
    o_k+c_k\mathbf e,&\mu_k=2,\\
    o_k,&\mu_k>2.
  \end{cases}
  \tag{3.6.45}\label{eq:3-6-45}
\end{equation}
注意，当 $\mu_k>0$ 时，这样的 $b_k(\mathrm{weak})$ 自动存在。（$o_k$
仍按定义 3.6.23 的同样方式定义。）

由于稍后将证明更一般的结果（定理 3.8.50），这里略去证明，留给读者。
当
\[
  (C,\mathfrak m)
  =\bigl(C(L;\Lambda_{0,\mathrm{nov}})^+,\mathfrak m^+\bigr),
\]
其中 $(C(L;\Lambda_{0,\mathrm{nov}}),\mathfrak m)$ 如第 3.5 节，而上标
$+$ 如第 3.3 节时，用 $o_k(L;\mathrm{weak})$ 代替
$o_k(\mathrm{weak})$。

\subsection{超势与 $\widehat{\mathcal M}(C)$}
\label{subsec:3-6-4}
\liftIndexRangeStart{I114}

回忆我们假设
$b\equiv0\pmod{\Lambda^+_{0,\mathrm{nov}}}$ 的理由：要使形式级数
$e^b=1+b+b\otimes b+\cdots$ 收敛。不过，如果存在
$k(\lambda)\in\mathbb Z$，满足
$\lim_{\lambda\to\infty}k(\lambda)=\infty$，并且只要
$k>k(\lambda)$ 就已经有
$\mathfrak m_k\equiv0\pmod{T^\lambda}$，那么对任意
$b\in C[1]^0$，$\widehat d(e^b)$ 都收敛
\translatornote{原书此处写 $b\in C^0[1]$；紧接的猜想 3.6.46 写作
$C(L)[1]^0$，而本节从定义 3.6.4 起始终要求 $b$ 的移位次数为零，故在此
恢复同一类型。}。切卡诺夫的工作 \textbf{[Chek02]} 中出现的
$A_\infty$ 代数确实如此。

在我们的情形中，这样的 $k(\lambda)$ 似乎并不存在（见下面的猜想
3.6.53）。不过，我们仍然期待如下结论。

\begin{conjecture}{3.6.46}\label{conj:3-6-46}
可以选取我们的滤过 $A_\infty$ 代数 $(C(L),\mathfrak m)$，使得对任意
$b\in C(L)[1]^0$，$\widehat d(e^b)$ 都收敛。
\end{conjecture}

可把这个猜想同引言中的猜想 T 比较。

\begin{remark}{3.6.47}\label{rem:3-6-47}
注意，猜想 3.6.46 意味着：对任意固定的 $\beta$ 和
$b\in C(L;\mathbb Q)$，级数
\[
  \sum_k\mathfrak m_{k,\beta}(b,\ldots,b)
\]
收敛。先验地说，得到这种收敛性应当比在算子 $\mathfrak m_k$ 中代入
$T=e^{-1}$ 后仍有收敛性更容易。也就是说，猜想 3.6.46 看起来远比证明
级数
\[
  \sum_{k,\beta}e^{-\omega[\beta]}
  \mathfrak m_{k,\beta}(b,\ldots,b)
\]
收敛容易；后一个级数需要对所有不同的 $\beta$ 求和。（另见
\textbf{[Fuk02I, KoSo01]} 等。）
\end{remark}

现在讨论：当 $L\subset M$ 是卡拉比--丘三维流形中的特殊拉格朗日子流形
（更一般地，是马斯洛夫类为零的拉格朗日子流形）时，Maurer--Cartan
方程预期具有怎样的性质。

首先，由断言滤过 $A_\infty$ 代数的典范模型存在的定理 5.4.2，可以取
\[
  C(L)=H(L;\mathbb Q)\otimes\Lambda_{0,\mathrm{nov}};
\]
见定义 5.4.3。我们还把 $\overline C(L)=H(L;\mathbb Q)$ 取作微分分次
代数。此外，可以略去零次与三次上同调（见第 5.4.6 小节），于是只剩
$H^1(L;\mathbb Q)$ 与 $H^2(L;\mathbb Q)$。

其次，有一个人们普遍相信为真但尚未证明的猜想。

\begin{conjecture}{3.6.48}\label{conj:3-6-48}
可以在 $(C(L),\mathfrak m)$ 上选取一个 $A_\infty$ 结构，使其满足下列
条件：存在一个由同调中的庞加莱对偶配对诱导的配对
\[
  \langle\ ,\ \rangle\colon
  C(L)\otimes C(L)\longrightarrow\Lambda_{0,\mathrm{nov}},
\]
它满足循环对称性条件\liftIndex{I113}
\[
\begin{aligned}
  \langle x_0,\mathfrak m_k(x_1,\ldots,x_k)\rangle
  ={}&(-1)^{\deg' x_k(\deg' x_0+\cdots+\deg' x_{k-1})}\\
  &\mathrel{}\cdot
  \langle x_k,\mathfrak m_k(x_0,x_1,\ldots,x_{k-1})\rangle.
\end{aligned}
\]
\end{conjecture}

如果可以忽略横截性和符号方面的微妙问题，猜想 3.6.48 很容易看出。
（在保持所需对称性的同时实现横截性，是一项困难的工作；不过可见第 A3
节和 \textbf{[Fuk05II]}。）在这些假设下，选取
$H^1(L;\mathbb Q)$ 的一组基 $e_i$（$i=1,\ldots,m$），并以
$b=\sum_i x_i e_i$ 定义函数
\[
  \Psi\colon
  \underbrace{\Lambda_{0,\mathrm{nov}}\times\cdots\times
  \Lambda_{0,\mathrm{nov}}}_{m\text{ 个}}
  \longrightarrow\Lambda_{0,\mathrm{nov}}
\]
如下：
\begin{equation}
  \qquad
  \Psi(x_1,\ldots,x_m)
  =\sum_k\frac{1}{k+1}
  \langle b,\mathfrak m_k(b,\ldots,b)\rangle.
  \tag{3.6.49}\label{eq:3-6-49}
\end{equation}
（见 \textbf{[Laza01, Tom01]}。）这个函数在物理学文献中似乎广为人知，
称为\emph{超势}。为了使~\eqref{eq:3-6-49} 的右端对
$b\in C(L;\Lambda_{0,\mathrm{nov}})$ 良定义，需要猜想 3.6.46。

\begin{proposition}{3.6.50}\label{prop:3-6-50}
假设猜想 3.6.46 和 3.6.48 成立。则 $b=\sum_i x_i e_i$ 满足
$\widehat d(e^b)=0$，当且仅当 $(x_1,\ldots,x_m)$ 是 $\Psi$ 的临界点。
\end{proposition}

\begin{proof}
利用猜想 3.6.48，得到
\[
  \frac{\partial}{\partial x_i}\Psi(x_1,\ldots,x_m)
  =\sum_k\langle e_i,\mathfrak m_k(b,\ldots,b)\rangle
  =\langle e_i,\mathfrak m(e^b)\rangle.
\]
命题随即成立。
\end{proof}

可以用类似方式加强这个命题。为此，定义
$\mathfrak m^b_{i_1,\ldots,i_k,i_0}\in\Lambda^+_{0,\mathrm{nov}}$ 与
$\mathfrak l^b_{i_1,\ldots,i_k,i_0}\in\Lambda^+_{0,\mathrm{nov}}$ 如下：
\begin{equation}
  \mathfrak m^b_{i_1,\ldots,i_k,i_0}
  =\bigl\langle e_{i_0},
    \mathfrak m(e^b e_{i_1}e^b\cdots e^b e_{i_k}e^b)
   \bigr\rangle,
  \tag{3.6.51.1}\label{eq:3-6-51-1}
\end{equation}
以及
\begin{equation}
\begin{aligned}
  \mathfrak l^b_{i_1,\ldots,i_k,i_0}
  &=\sum_{\sigma\in\mathfrak S_k}
    \mathfrak m^b_{i_{\sigma(1)},\ldots,i_{\sigma(k)},i_0}\\
  &=\frac{1}{k+1}\sum_{\sigma\in\mathfrak S_{k+1}}
    \mathfrak m^b_{i_{\sigma(1)},\ldots,i_{\sigma(k)},i_{\sigma(0)}}.
\end{aligned}
\tag{3.6.51.2}\label{eq:3-6-51-2}
\end{equation}
这里 $\mathfrak S_k$ 是 $k$ 个字母的对称群，其阶为 $k!$。于是
$\mathfrak l^b_{i_1,\ldots,i_k,i_0}$ 是由 $\mathfrak m$ 对称化所得
$L_\infty$ 结构的结构常数（见第 A3 节）。

\begin{proposition}{3.6.52}\label{prop:3-6-52}
假设猜想 3.6.46 和 3.6.48 成立。则
\[
  \frac{\partial^{k+1}\Psi}
  {\partial x_{i_1}\cdots\partial x_{i_k}\partial x_{i_0}}(b)
  =\mathfrak l^b_{i_1,\ldots,i_k,i_0}.
\]
\end{proposition}

\begin{proof}
若 $b_0=\sum_i x_i^0e_i$、$b=\sum_i x_i e_i$，则
\[
\begin{aligned}
&\Psi(x_1+x_1^0,\ldots,x_m+x_m^0)\\
&=\sum_k\frac{1}{k+1}
  \langle b_0+b,\mathfrak m_k(b_0+b,\ldots,b_0+b)\rangle\\
&=\sum_{k,\ell_0,\ldots,\ell_k}
  \frac{1}{k+1+\sum_i\ell_i}
  \Bigl\langle b,
  \mathfrak m_{k+\sum_i\ell_i}
  (b_0^{\ell_0},b,b_0^{\ell_1},\ldots,
   b_0^{\ell_{k-1}},b,b_0^{\ell_k})\Bigr\rangle\\
&\quad+
  \sum_{k,\ell_0,\ldots,\ell_k}
  \frac{1}{k+1+\sum_i\ell_i}
  \Bigl\langle b_0,
  \mathfrak m_{k+\sum_i\ell_i}
  (b_0^{\ell_0},b,b_0^{\ell_1},\ldots,
   b_0^{\ell_{k-1}},b,b_0^{\ell_k})\Bigr\rangle\\
&\mathrel{\overset{*}{=}}
  \sum_{k,\ell_0,\ldots,\ell_k}
  \frac{1}{k+1}
  \Bigl\langle b,
  \mathfrak m_{k+\sum_i\ell_i}
  (b_0^{\ell_0},b,b_0^{\ell_1},\ldots,
   b_0^{\ell_{k-1}},b,b_0^{\ell_k})\Bigr\rangle\\
&=\sum_k\frac{1}{k+1}
  \langle b_0,\mathfrak m_k(b_0,\ldots,b_0)\rangle
  +\sum_k\frac{1}{k+1}
  \langle b,\mathfrak m_k^{b_0}(b,\ldots,b)\rangle.
\end{aligned}
\]
对于上面的等号 $\overset{*}{=}$，循环对称性（猜想 3.6.48）的用法如下。
对 $k$ 与 $\ell$，考虑和
\[
  \sum_{\ell_0+\ell_1+\cdots+\ell_k=\ell}
  \Bigl\langle b,
  \mathfrak m_{k+\ell}
  (b_0^{\ell_0},b,b_0^{\ell_1},\ldots,
   b_0^{\ell_{k-1}},b,b_0^{\ell_k})\Bigr\rangle.
\]
对其中的变量作循环置换。在 $\ell+k+1$ 个循环置换中，有 $k+1$ 个具有
下列形式：
\[
  \sum_{\ell_0+\ell_1+\cdots+\ell_k=\ell}
  \Bigl\langle b,
  \mathfrak m_{k+\ell}
  (b_0^{\ell_0},b,b_0^{\ell_1},\ldots,
   b_0^{\ell_{k-1}},b,b_0^{\ell_k})\Bigr\rangle;
\]
另有 $\ell$ 个具有下列形式：
\[
  \sum_{\ell_0+\ell_1+\cdots+\ell_{k+1}=\ell-1}
  \Bigl\langle b_0,
  \mathfrak m_{k+\ell}
  (b_0^{\ell_0},b,b_0^{\ell_1},\ldots,
   b_0^{\ell_k},b,b_0^{\ell_{k+1}})\Bigr\rangle.
\]
这就证明了等号 $\overset{*}{=}$。

因此，把 $\mathfrak m_k$ 换成 $\mathfrak m_k^{b_0}$ 后，只需在 $b=0$
处证明命题。此时
\[
\begin{aligned}
  &\frac{\partial^{k+1}\Psi}
  {\partial x_{i_1}\cdots\partial x_{i_k}\partial x_{i_0}}(0)\\
  &\quad={}
  \left.
  \frac{1}{k+1}
  \frac{\partial^{k+1}}
  {\partial x_{i_1}\cdots\partial x_{i_k}\partial x_{i_0}}
  \langle b,\mathfrak m_k(b,\ldots,b)\rangle
  \right|_{b=0}
  =\mathfrak l_{i_1,\ldots,i_k,i_0},
\end{aligned}
\]
正是所需结论。
\end{proof}

现在考虑卡拉比--丘三维流形 $M$ 及其马斯洛夫类为零的拉格朗日子流形
$L\subset M$ 所成的对子 $(M,L)$。选取适当的殆复结构 $J$ 后，可以假设
所有在某处单射的全纯圆盘都是孤立的。把它们枚举为
\[
  u_i\colon(D^2,\partial D^2)\longrightarrow(M,L),
\]
并记为下式。\translatornote{原书第二式写
$[\partial u_i]\cap x_j$，但前文已定义 $e_j$ 为
$H^1(L;\mathbb Q)$ 的基、$x_j$ 为其标量坐标；边界同调类只能与前者
配对，故据类型关系改为 $e_j$。}
\[
  \lambda_i=\int_{D^2}u_i^*\omega,
  \qquad
  a_{i,j}=[\partial u_i]\cap e_j.
\]

\begin{conjecture}{3.6.53}\label{conj:3-6-53}
设 $\Psi$ 是上面定义的超势。则
\[
  \Psi(x_1,\ldots,x_m)
  =c+\sum_{i=1}^{\infty}\sum_{d=1}^{\infty}
  \epsilon_{i,d}n_dT^{d\lambda_i}
  \exp\left(d\sum_{j=1}^m a_{i,j}x_j\right).
\]
\end{conjecture}

这里 $\exp\colon\mathbb C\to\mathbb C$ 是通常的指数函数；$n_d$ 是只依赖
$d$ 的数，用来描述如何计数圆盘 $u_i$ 的 $d$ 重覆盖。（一些文献
\textbf{[OoVa96, KatLi01]} 中的结果暗示 $n_d=d^{-2}$。）
$\epsilon_{i,d}=\pm1$ 是计数圆盘 $u_i$ 的 $d$ 重覆盖时所用的符号。
$c\in\Lambda_{0,\mathrm{nov}}$ 与 $(x_1,\ldots,x_m)$ 无关，但可能依赖于
$L$、$J$ 和扰动。

作者预计，可在不久的将来利用 \textbf{[Fuk05II]} 的技术证明本小节所述的
若干猜想。

\begin{remark}{3.6.54}\label{rem:3-6-54}
在环面情形中，超势与弗洛尔上同调之间的关系见
\textbf{[ChOh03, Cho04II, FOOO08I]}。Cheol-Hyun Cho 在
\textbf{[Cho04II]} 中证明了一个与命题 3.6.52 类似的公式。在环面情形中，
令 $T=\exp(-1)$，就可以把朗道--金兹堡势与我们的势函数等同起来。

在卡拉比--丘情形中，超势与滤过 $A_\infty$ 结构之间的关系同环面情形
具有不同形式：在本节所讨论的卡拉比--丘情形中，超势的一阶导数与
$\mathfrak m_0(1)$ 有关；而超势的 $k$ 阶导数与
$\mathfrak m_{k-1}$ 的对称化有关。在环面情形中，$\mathfrak m_0(1)$
本身就是朗道--金兹堡势函数，它的一阶导数对应于 $\mathfrak m_1$。
\liftIndexRangeEnd{I114}
\end{remark}
